\documentclass[conference]{IEEEtran}
\usepackage[utf8]{inputenc}
\usepackage[T1]{fontenc}
\usepackage{cite}
\usepackage{amsmath,amssymb}
\usepackage{graphicx}
\usepackage{booktabs}
\usepackage{url}

\title{Generate--Verify--Repair for Intent‑Driven Network Changes: Controlled Evaluation with a Deterministic Verifier}

\author{
\IEEEauthorblockN{Autonomous AI Researcher}
\IEEEauthorblockA{CNSM 2026 Agentic AI Researcher Track}
}

\begin{document}

\maketitle

\begin{abstract}
We evaluate whether a generate--verify--repair (GVR) pipeline---single‑shot LLM generation followed by a deterministic network verifier and up to one automated repair---reduces accepted unsafe network changes versus LLM‑only generation on a standardized synthetic NetOps intent workload. In a preregistered paired design (study-hosted-netops-gvr-v1-2026-08-29) we ran N=500 distinct synthetic intents (shared single initial candidate per intent) with generation by gpt-5-mini and adjudication by a deterministic verifier. Results: guarded (GVR) accepted 500/500 final changes; baseline (LLM-only) accepted 496/500 (paired difference = 0.008). The preregistered primary exact conditional McNemar two‑sided test on the discordant pairs (n10 = 4, n01 = 0) yields p = 0.125 (computed as p = 2 * PBinom(X >= 4; n = 4, p = 0.5)). An exploratory paired bootstrap percentile 95\% CI (10,000 resamples, seed = 1729) = [0.002, 0.016] (bootstrap labeled exploratory per preregistration and interpreted cautiously given only four discordants). Repairs were invoked for 4 initial candidates; total model calls used = 504 (500 shared\_initial + 4 repair calls). The preregistered templated‑automation comparator was not executed. Because the preregistered decision rule is the exact conditional McNemar (p = 0.125), we do not claim confirmatory statistical significance. We reconcile the conditional McNemar and the exploratory bootstrap CI, document verifier scope and known blind spots, provide execution accounting and representative validator traces, and give explicit artifact pointers and hashes for reproducibility.
\end{abstract}

\section{Introduction}
Operator intent\ensuremath{\rightarrow}device configuration translation can reduce manual effort, but errors in generated CLI sequences or transient unsafe states can cause outages. Modern large language models (LLMs) show promise for translating human intent to device configuration, yet hallucination and factual errors remain central reliability concerns for operational NetOps. A practical mitigation architecture is generate--verify--repair (GVR): produce a candidate configuration with an LLM, deterministically verify it against encoded workflow and policy rules, and, if verification fails, run a bounded automated repair attempt. This paper reports a preregistered, artifact‑grounded paired experiment that asks: Does a GVR pipeline (LLM \ensuremath{\rightarrow} deterministic verifier \ensuremath{\rightarrow} up to one automated repair) produce fewer accepted unsafe changes than LLM‑only generation on a standardized synthetic NetOps intent workload? We contribute: (1) a preregistered paired trial (N=500 paired intents) executed end‑to‑end with archived artifacts; (2) an artifact‑grounded description of generation, verifier semantics, and the bounded repair loop; (3) reconciled confirmatory inference (two‑sided conditional exact McNemar) and exploratory paired bootstrap CI descriptions with an explicit small‑discordant‑count caveat; and (4) execution accounting, representative validator traces, and a discussion of verifier blind spots and practical external‑validity limits.

\section{Related Work}
Three converging literatures inform the reliability challenge for LLM‑driven NetOps and motivate a GVR architecture. First, survey and systems work documents active interest in using LLMs for intent translation, configuration generation, AIOps, and multi‑agent orchestration, while emphasizing hallucination and factual errors as deployment barriers. Large surveys and system descriptions synthesize these trends and identify practical obstacles such as vendor heterogeneity, control‑plane timing, and dataset availability [doi:10.1145/3718958.3750537; https://openalex.org/W4411735779]. Intent‑driven configuration prototypes and frameworks further demonstrate feasibility of LLM‑based translation but typically pair those generators with ad‑hoc validation or operator oversight rather than large, controlled trials against deterministic verifiers [doi:10.1002/nem.2313]. These works establish that translation quality alone is insufficient for operational acceptance and motivate tool‑augmented pipelines that embed validation as a gate.

Second, methodological and tool‑integration work advocates for explicit generate‑validate‑repair (or RAG/tool‑augmented) evaluation practices and standardized validator semantics. Recent methodological papers propose blueprints for evaluating retrieval‑ or tool‑augmented LLMs in operational contexts and call for measuring verifier coverage, failure modes, and orchestration costs rather than only generator fidelity [https://openalex.org/W4403443637]. These studies highlight three evaluation gaps that our experiment targets: (a) need for end‑to‑end trial designs that preserve pairing between conditions on identical intents; (b) the value of frozen preregistration to avoid post‑hoc analysis choices; and (c) the importance of quantifying model‑call overhead per avoided unsafe change when verifiers are invoked.

Third, deterministic configuration verifiers and scalable verification systems provide concrete oracles for many syntactic and workflow properties but exhibit known modeling limits. Work on scalable, distributed verifiers (Batfish derivatives and localized subspecification approaches) demonstrates that automated checks can detect availability/sequence and policy violations at scale, yet these verifiers vary in which control‑plane or vendor semantics they model and therefore differ in the set of hazards they can detect [https://openalex.org/W4413757123]. Prior verifier papers typically measure verifier expressiveness and scalability; they do not, however, couple that verifier deterministically with an LLM generator in a preregistered paired trial that then measures repaired vs accepted unsafe changes across a broad task bank. Our study uses a deterministic\_netops\_validator\_v5 that implements canonical parsing, required‑sequence/workflow checks, availability/group constraints, and paired‑field dependency rules---i.e., the class of properties many verifier systems aim to cover---but we acknowledge, as prior work does, that verifier pass does not imply absence of vendor‑specific or timing hazards.

Across these literatures, two important empirical gaps remain. First, although several system papers and frameworks argue for verifier‑assisted or agentic LLM pipelines, there are few publicly archived, preregistered paired comparisons that measure how bounded repair loops change the rate of verifier‑detectable unsafe outputs across a large, stratified set of intents; most prior evaluations are descriptive, single‑condition, or lack frozen analysis plans [https://openalex.org/W4413757586]. Second, methodological reviews emphasize that evaluation must report verifier codebooks, failure‑mode counts, and orchestration costs (model calls, repair rates) to allow operational judgment [https://openalex.org/W4403443637]; without those details, claims about reliability gains remain qualitative. Our experiment addresses both gaps by preregistering a paired confirmatory estimand, archiving the verifier codebook and representative validator traces, reporting provider audit counts (model\_calls\_used = 504; repair = 4), and publishing the paired contingency table for reproducible exact‑test calculation. In short, we position our work as an artifact‑grounded trial that operationalizes the recommendations from recent methodological work and connects generator performance to deterministic verifier coverage in a controlled paired design.

\section{Methodology}
Preregistration and confirmatory plan: The experimental design, sampling, estimands, and primary decision rule were preregistered as study-hosted-netops-gvr-v1-2026-08-29. The primary estimand is the paired success‑rate difference (GVR - baseline) across N = 500 distinct intents stratified evenly over 10 synthetic topology profiles (50 intents per profile). Sampling seeds, the deterministic task generator, and the analysis executor were frozen prior to execution (preregistration manifest SHA‑256 = aa8982a2\allowbreak{}7c73e515\allowbreak{}21ac023a\allowbreak{}78adcfa3\allowbreak{}58ef3978\allowbreak{}c0f9bc05\allowbreak{}213801fc\allowbreak{}bd2f1d0a). The preregistered confirmatory decision rule is the two‑sided conditional exact McNemar computed on the observed discordant pairs; the preregistration also specified an exploratory paired bootstrap percentile interval (10,000 resamples; bootstrap\_seed = 1729) as a descriptive interval only.

Task generation and paired design: A deterministic task generator (deterministic\_netops\_task\_generator\_v5) produced 500 distinct intent tasks (task\_manifest entries archived). Each intent was executed under two paired conditions using the same shared\_initial\_generation candidate: (1) baseline---accept the shared initial LLM candidate as the final change; and (2) guarded (GVR)---validate the shared candidate with the deterministic verifier and, if failing, invoke at most one automated repair attempt. Pairing ensured within‑intent control of task difficulty, topology profile, and the single shared initial candidate design as preregistered.

Generation model and orchestration: The declared generation model for the run is gpt-5-mini (declared\_model\_version recorded in execution/model\_configuration.json; execution/model\_configuration.json SHA‑256 = a40e96e2\allowbreak{}12ab87da\allowbreak{}251ac0d6\allowbreak{}dbb75bc5\allowbreak{}43afa134\allowbreak{}38b5a6b9\allowbreak{}ebd07359\allowbreak{}7ffdd820). Per the execution manifest and provider audit, the run produced 500 shared\_initial\_generation calls and 4 repair calls (hosted\_model\_call\_event\_count = 504; model\_calls\_used = 504). The orchestration adapter family is hosted\_netops\_gvr\_v1 using the hosted provider recorded as 'openai\_responses'. The maximum\_repair\_calls\_per\_task was set to 1 in the preregistration and enforced by the controller.

Deterministic verifier: role, semantics, and outputs: The primary adjudicator is deterministic\_netops\_validator\_v5 (validator\_version = 5.0), implemented as a deterministic canonicalization and rule‑checking pipeline. For each candidate CLI sequence, the verifier executes the following stages in order: (1) syntactic parsing and canonical normalization (tokenization of interface commands, canonical ordering of multi‑line sequences, removal of formatting variants); (2) command‑level required‑sequence enforcement (patterns such as must\_be\_down\_for\_fields, all\_down\_before\_any\_field\_change, make\_before\_break switchover); (3) availability/group constraints (exactly\_active, minimum\_active\_in\_groups); (4) dependency and paired‑field constraints (paired MTU changes, equal\_final\_fields between interfaces); and (5) final‑state comparison against the task's expected\_final\_state. Any nonempty violation set yields a failing outcome. For each episode the verifier emits a structured validator\_trace containing fields including normalized\_configuration, observed\_touched\_field\_count, parsed\_command\_count, required\_command\_count, violation\_count and list, validation\_before/after summaries, and boolean valid. Representative verifier\_trace objects are archived under execution/scoring/ (see representative artifact examples). Passing the verifier is a necessary experimental safety gate but is not claimed to be sufficient for all real‑world hazards (see Threats to Validity).

Verifier codebook and artifact pointer: The human‑readable verifier codebook (violation‑code \ensuremath{\rightarrow} short description) used to construct repair prompts and to interpret validator\_trace violations is archived as an artifact referenced by the execution manifest. The execution bundle includes a full hash manifest; the execution manifest path is execution/execution\_manifest.json (the manifest enumerates the verifier codebook artifact path and its SHA‑256 for direct retrieval). Representative scoring JSON objects in execution/scoring/ contain the violation codes emitted for each failing candidate.

Repair policy, prompt template, and acceptance rule: When the verifier returns valid=false for a shared initial candidate under the guarded condition, the GVR controller issues at most one repair prompt to the same generation model. The repair prompt is constrained and intentionally compact: it contains (a) the original high‑level intent abstraction for the task, (b) the verifier violation codes and short descriptions extracted from the verifier\_trace, and (c) a request for a corrected canonical CLI sequence that satisfies the listed violation codes. The repair template minimizes additional contextual material to isolate the marginal effect of one repair attempt. A repair candidate is accepted only if a subsequent verifier invocation returns valid=true. All repair prompts, responses, and verifier outcomes are logged and archived in execution/provider\_calls/ and execution/scoring/.

Scoring rule and outcome definition: Per preregistration, each intent's final accepted change is scored binary by the primary verifier: success = final accepted configuration passes deterministic\_netops\_validator\_v5 (valid=true) versus unsafe = final accepted configuration fails (valid=false). For baseline pairs the shared candidate is the accepted change without verification‑based repair; for guarded pairs acceptance required verifier pass after repair (if any). The analysis used complete pairs only (complete\_pair\_count = 500) per the preregistered missingness plan; technical failures are recorded and were zero in this run (failed\_episode\_count = 0).

Statistical procedures and reproducible calculations: The preregistered primary test is the two‑sided conditional exact McNemar operating on discordant pairs. Let n10 be the number of pairs where guarded succeeds and baseline fails, and n01 the reverse; n = n10 + n01 and k = n10. The conditional exact two‑sided p‑value is computed as p = 2 * PBinom(X >= k; n, 0.5), with clipping at 1.0. In this run the archived paired contingency (analysis/paired\_contingency\_table.csv; archive hash listed in the evidence bundle) is n11 = 496, n10 = 4, n01 = 0, n00 = 0, so n = 4 and k = 4. For n = 4, PBinom(X >= 4; n = 4, p = 0.5) = (1/16), giving p = 2*(1/16) = 0.125. The exact McNemar implementation and this conditional binomial formula were preregistered and applied exactly as documented here; the contingency CSV and analysis logs are archived for reproducibility.

Exploratory paired bootstrap: The preregistration additionally specified an exploratory paired bootstrap percentile CI for descriptive interpretation. We resampled paired outcomes with replacement (10,000 resamples; bootstrap\_seed = 1729) and computed the paired success‑rate difference (guarded - baseline) per resample; the 2.5th and 97.5th percentiles produce the reported exploratory interval [0.002, 0.016]. The bootstrap was explicitly labeled exploratory in the preregistration because percentile intervals can have poor frequentist coverage when discordant counts are very small and sample margins are near ceiling; we therefore present the bootstrap as a descriptive interval and rely on the preregistered exact McNemar for the primary confirmatory decision.

Provider accounting, determinism, and retry policy: Provider audit fields and the deterministic reconciliation artifact record hosted\_model\_call\_event\_count = 504 (completed\_hosted\_model\_call\_event\_count = 504; failed\_hosted\_model\_call\_event\_count = 0), stage\_counts = \{shared\_initial\_generation: 500, repair: 4\}, and model\_calls\_used = 504. The controller enforced maximum\_attempts\_per\_call = 1 and maximum\_repair\_calls\_per\_task = 1 as preregistered. All random seeds used for task generation, bootstrap resampling, and other nondeterministic steps are recorded in the archived artifacts to ensure computational reproducibility.

Predefined exclusion, missingness, and sensitivity: Per preregistration there were no exclusions applied; complete\_pair\_count = 500 and no technical failures occurred. The preregistration specified a sensitivity analysis treating technical failures as unsafe; because no such failures occurred the primary and sensitivity treatments coincide. All missingness, contamination, and reconciliation artifacts are archived (analysis/missingness\_summary.csv, analysis/contamination\_summary.csv, deterministic\_reconciliation.json).

Reproducibility artifacts and pointers: Key archived artifacts referenced in this methodology include execution/raw\_results.jsonl (SHA‑256 = fe9b93ef\allowbreak{}5cc5a237\allowbreak{}a3361f18\allowbreak{}f45fabf1\allowbreak{}4cd23a1c\allowbreak{}679556f6\allowbreak{}5477f9f6\allowbreak{}81915d02), execution/task\_manifest.jsonl (SHA‑256 = 5a822e78\allowbreak{}7302a0b3\allowbreak{}bb03a86a\allowbreak{}81b20564\allowbreak{}a44e2636\allowbreak{}731f853f\allowbreak{}9d45ac12\allowbreak{}e4876cc8), analysis/paired\_contingency\_table.csv (SHA‑256 = 8fb135ce\allowbreak{}c541cbe0\allowbreak{}b14f1612\allowbreak{}5db33154\allowbreak{}a41460c4\allowbreak{}9958427f\allowbreak{}e8e7f3e1\allowbreak{}f2abfece), and analysis/analysis\_log.jsonl (SHA‑256 = 261ec80e\allowbreak{}a9343fcc\allowbreak{}bcd0a622\allowbreak{}94f53827\allowbreak{}73770fd5\allowbreak{}9d387dc9\allowbreak{}d749de50\allowbreak{}ca909b3e). The full execution manifest (execution/execution\_manifest.json) contains the complete list of artifact paths and SHA‑256 digests for direct retrieval. Reproducible calculation parameters (bootstrap resamples, bootstrap\_seed = 1729) and the exact McNemar formula are recorded in analysis/analysis\_log.jsonl.

Implementation notes and limits applied to the methodology: The verifier enforces canonical, task‑specified properties and exact final‑state equality against the synthetic expected\_final\_state; it intentionally omits vendor‑specific parser idiosyncrasies, control‑plane timing/simulation, and out‑of‑band hardware constraints. The repair policy is intentionally bounded (single repair) to limit model‑call overhead and to measure the marginal corrective value of one repair attempt in an operationally conservative budget. These design choices were preregistered and are explicitly documented to bound claims about operational generality.

\section{Results}
Execution and primary counts: The run completed completed\_episode\_count = 1000 (500 intents \ensuremath{\times} 2 conditions) with model\_calls\_used = 504. analysis/condition\_summary.csv (archived hash: f3e19742\allowbreak{}5cc03dec\allowbreak{}41cda77f\allowbreak{}078f6d0b\allowbreak{}079b06bb\allowbreak{}de7ac736\allowbreak{}de5fe988\allowbreak{}b027ac58) reports baseline successes 496/500 (0.992) and guarded successes 500/500 (1.000). The paired contingency table (analysis/paired\_contingency\_table.csv; archived hash: 8fb135ce\allowbreak{}c541cbe0\allowbreak{}b14f1612\allowbreak{}5db33154\allowbreak{}a41460c4\allowbreak{}9958427f\allowbreak{}e8e7f3e1\allowbreak{}f2abfece) is: n11 = 496, n10 = 4, n01 = 0, n00 = 0. Repairs were invoked for four initial candidates; all four repair attempts resulted in subsequent verifier passes and correspond to the four discordant pairs (n10 = 4). Mean model calls per accepted change \ensuremath{\approx} 504/500 = 1.008.

Exact McNemar implementation and numeric reproducibility: To remove ambiguity we state the exact McNemar implementation used and the input counts. The preregistered conditional exact McNemar test we applied conditions on the observed discordant total n = n10 + n01 and uses the two‑sided binomial tail on k = n10 with success probability 0.5. For our observed contingency (n10 = 4, n01 = 0) we have n = 4 and k = 4. The one‑sided binomial tail PBinom(X >= 4; n = 4, p = 0.5) = (1/2)\textasciicircum{}4 = 1/16; the two‑sided p is computed as p = 2 * PBinom(X >= k; n, 0.5) = 2 * (1/16) = 0.125 (clipped to <=1.0). This precise computation reproduces the reported preregistered p‑value and is directly reproducible from the archived paired\_contingency\_table.csv artifact (hash above).

Exploratory paired bootstrap interval and interpretation: The preregistered exploratory paired bootstrap percentile 95\% CI (10,000 resamples; seed = 1729; artifacts archived as analysis/paired\_difference\_figure.svg and analysis/analysis\_log.jsonl; analysis\_log.jsonl SHA‑256 = 261ec80e\allowbreak{}a9343fcc\allowbreak{}bcd0a622\allowbreak{}94f53827\allowbreak{}73770fd5\allowbreak{}9d387dc9\allowbreak{}d749de50\allowbreak{}ca909b3e) for the paired difference is [0.002, 0.016]. We emphasize and clarify two points requested by reviewers: (1) this bootstrap CI was predeclared exploratory (not the preregistered primary inferential basis), and (2) the bootstrap mechanism differs from the conditional McNemar in an important conditioning property that explains the apparent discrepancy when discordant counts are very small. Concretely, the exact conditional McNemar conditions on the observed discordant margin n and evaluates the binomial tail of k under n with p = 0.5; with n = 4 this yields a discrete two‑sided p = 0.125. The paired bootstrap instead resamples entire paired outcomes (not conditioning on fixed discordant margins), so percentile intervals can exclude zero when almost all pairs are concordant and a small number of discordant pairs consistently favor one direction across resamples. Because the bootstrap percentile interval has known coverage limitations when discordant counts are tiny and sampling variability of the margins is large relative to n, we label the bootstrap interval explicitly exploratory and advise cautious interpretation; the preregistered decision rule remains the exact conditional McNemar p = 0.125.

Representative artifact diagnostics and codebook retrieval: Representative scoring JSON objects (execution/scoring/task-00000*-*.json) show verifier\_trace entries with validation\_before, validation\_after, normalized\_configuration, parsed\_command\_count, required\_command\_count, violation\_count, and validator\_version = 5.0. The four repaired episodes show verifier‑detectable sequence or availability violations that a single constrained repair prompt corrected; repair prompts included the verifier violation codes and the intent abstraction per the preregistered policy. The human‑readable verifier codebook (violation‑code \ensuremath{\rightarrow} description) is archived and its artifact path entry is enumerated in the archived execution manifest (execution/execution\_manifest.json) for direct retrieval; the execution manifest is part of the evidence bundle and lists the exact codebook artifact path and SHA‑256 so readers can obtain the precise codebook used for scoring. The paired contingency CSV hash is analysis/paired\_contingency\_table.csv (SHA‑256 = 8fb135ce\allowbreak{}c541cbe0\allowbreak{}b14f1612\allowbreak{}5db33154\allowbreak{}a41460c4\allowbreak{}9958427f\allowbreak{}e8e7f3e1\allowbreak{}f2abfece) and the raw results JSONL hash is execution/raw\_results.jsonl (SHA‑256 = fe9b93ef\allowbreak{}5cc5a237\allowbreak{}a3361f18\allowbreak{}f45fabf1\allowbreak{}4cd23a1c\allowbreak{}679556f6\allowbreak{}5477f9f6\allowbreak{}81915d02).

Provider audit and execution manifest: Deterministic\_reconciliation.json records provider\_call\_audit: hosted\_model\_call\_event\_count = 504, completed\_hosted\_model\_call\_event\_count = 504, failed\_hosted\_model\_call\_event\_count = 0, cache\_hit\_count = 0, cache\_miss\_count = 504, and stage\_counts: shared\_initial\_generation = 500 and repair = 4. Representative raw artifacts and hashes cited here include execution/raw\_results.jsonl (SHA‑256 = fe9b93ef\allowbreak{}5cc5a237\allowbreak{}a3361f18\allowbreak{}f45fabf1\allowbreak{}4cd23a1c\allowbreak{}679556f6\allowbreak{}5477f9f6\allowbreak{}81915d02) and execution/task\_manifest.jsonl (SHA‑256 = 5a822e78\allowbreak{}7302a0b3\allowbreak{}bb03a86a\allowbreak{}81b20564\allowbreak{}a44e2636\allowbreak{}731f853f\allowbreak{}9d45ac12\allowbreak{}e4876cc8). All artifacts cited are retrievable from the archived execution manifest in the evidence bundle.

\section{Discussion}
Interpretation and reconciliation of inferential outcomes: The GVR pipeline prevented four verifier‑detectable unsafe initial candidates that baseline would have accepted. Under the preregistered conditional exact McNemar (p = 0.125) this difference is not confirmatory. The exploratory bootstrap CI excludes zero but should be interpreted cautiously: the conditional McNemar conditions on observed margins and evaluates the binomial tail of discordant outcomes; with n = 4 and k = 4 this yields p = 0.125. The bootstrap resamples pairs without fixing margins and can produce a nonzero lower CI bound when most pairs are concordant and a small number are discordant. Because the bootstrap was predeclared exploratory and because small discordant counts reduce frequentist coverage guarantees for percentile intervals, we report both inferential perspectives and follow the preregistered decision rule for the primary claim.

Verifier coverage and concrete blind spots: The deterministic verifier enforces canonical normalization, required sequencing, availability/group constraints, dependency and paired‑field rules, and exact final‑state equality to the task expected\_final\_state. It does not model vendor‑specific CLI semantics, timing/convergence effects in control‑plane protocols, vendor idiosyncrasies, hardware resource constraints, or out‑of‑band side effects. Passing this verifier is necessary for experiment safety but not sufficient for all production hazards; expanding verifier fidelity (vendor parsers, control‑plane simulators, timing models) is required before operational automation. The human‑readable verifier codebook (violation‑code \ensuremath{\rightarrow} description) is enumerated in the execution manifest (see execution/execution\_manifest.json in the evidence bundle) and the manifest lists the exact codebook artifact path and its SHA‑256 for direct retrieval.

Operational implications and cost tradeoffs: A compact deterministic verifier plus a bounded repair loop intercepted verifier‑detectable hazards at low LLM overhead in this synthetic workload. Observed repair rate was 4/500 (0.8\%), implying modest additional model calls under the bounded policy. Production adoption requires (a) measuring verifier runtime and scalability on real inventories, (b) extending verifier fidelity to vendor/timing effects, and (c) evaluating adversarial inputs where repair frequency and cost may increase. The provider audit (hosted\_model\_call\_event\_count = 504) quantifies the empirical model‑call overhead for this run.

Threats to validity and generality: Internal validity is strong because of preregistration, fixed seeds, and complete pairs (complete\_pair\_count = 500; execution\_manifest failed\_episode\_count = 0). Construct validity is limited by verifier abstractions and synthetic device templates (synthetic\_topologies\_v1 and synthetic\_device\_configs\_v1). External validity is constrained by single‑model scope (gpt-5-mini) and synthetic topologies; multi‑model replication and real vendor CLI tests are required to generalize. Conclusion validity was affected by realized margins: the preregistered power plan targeted an absolute paired increase \ensuremath{\approx} 0.08; the observed near‑ceiling baseline (0.992) produced only four discordant pairs and reduced realized power relative to the preregistered assumption. These limitations are documented in the archived manifest and preregistration.

\section{Limitations}
\begin{itemize}
\item Synthetic tasks and topology profiles were used (synthetic\_topologies\_v1 and synthetic\_device\_configs\_v1); external validity to production hardware, vendor‑specific CLIs, live telemetry, and control‑plane timing effects is untested.
\item Only a single generation model (gpt-5-mini) was evaluated; multi‑model generality and replication remain to be established.
\item Safety in this experiment is defined relative to deterministic\_netops\_validator\_v5's encoded rule set (syntactic parsing/normalization, required‑sequence/workflow‑policy checks, protected‑interface rules, group and availability constraints, dependency‑rule enforcement, and final‑state comparison to the task expected\_final\_state). Passing this verifier is necessary for the experiment's outcome but not sufficient for all real‑world safety properties.
\item The preregistered power plan targeted an absolute paired increase \ensuremath{\approx} 0.08 (8 percentage points). The observed baseline success rate (496/500 = 0.992) produced only four discordant pairs, substantially reducing realized power relative to the preregistered assumption and limiting confirmatory sensitivity.
\item The preregistered templated‑automation comparator (secondary confirmatory arm) was not executed; therefore the preregistered multiplicity plan (two confirmatory tests with Bonferroni correction) could not be fully applied and the familywise error interpretation is partial.
\end{itemize}

\section{Conclusion}
In a preregistered paired trial (N = 500) using gpt-5-mini and a deterministic verifier, a generate--verify--repair pipeline repaired four verifier‑detectable unsafe initial candidates at modest model‑call cost (4 repairs; model\_calls\_used = 504). The paired difference = 0.008 yields an exploratory bootstrap 95\% CI [0.002, 0.016], while the preregistered exact conditional McNemar test (n10 = 4, n01 = 0) yields p = 0.125; under the preregistered decision rule we do not claim confirmatory significance. Deterministic verification plus a bounded repair loop can remove verifier‑detectable unsafe outputs with low overhead in synthetic settings; broader operational claims require expanded verifier fidelity, adversarial NetOps evaluation, multi‑model replication, and execution of the preregistered templated comparator (which was not executed in this run). All primary execution and analysis artifacts cited here are archived in the evidence bundle for independent verification; see the archived execution manifest for full artifact paths and hashes.

\section*{Disclosure Statement}
This study was executed autonomously after an explicit autonomy lock. Generation model used in the archived run: gpt-5-mini (declared\_model\_version recorded in execution/model\_configuration.json; execution/model\_configuration.json SHA‑256 = a40e96e2\allowbreak{}12ab87da\allowbreak{}251ac0d6\allowbreak{}dbb75bc5\allowbreak{}43afa134\allowbreak{}38b5a6b9\allowbreak{}ebd07359\allowbreak{}7ffdd820). Orchestration adapter family: hosted\_netops\_gvr\_v1 using the hosted provider recorded as 'openai\_responses'. Key automated components recorded in the run: deterministic\_netops\_task\_generator\_v5 (task synthesis), deterministic\_netops\_validator\_v5 (primary deterministic verifier/scorer), and the GVR controller implementing shared\_initial\_generation plus an automated single‑repair step (maximum\_repair\_calls\_per\_task = 1). Preregistration: study-hosted-netops-gvr-v1-2026-08-29 (preregistration SHA‑256 = aa8982a2\allowbreak{}7c73e515\allowbreak{}21ac023a\allowbreak{}78adcfa3\allowbreak{}58ef3978\allowbreak{}c0f9bc05\allowbreak{}213801fc\allowbreak{}bd2f1d0a). The immutable initial master prompt is archived at provenance/master\_prompt.txt (SHA‑256 = 1872df1e\allowbreak{}1805d2d9\allowbreak{}6940456c\allowbreak{}a016bd66\allowbreak{}5d1d5196\allowbreak{}add77f5a\allowbreak{}cdf1582b\allowbreak{}b39b15ba). Primary high‑level execution quantities reported here (N = 500 complete pairs; model\_calls\_used = 504; repair calls = 4; paired counts n11 = 496, n10 = 4, n01 = 0, n00 = 0) are derived from archived execution and analysis artifacts. Representative artifact hashes included in the evidence bundle: execution/raw\_results.jsonl SHA‑256 = fe9b93ef\allowbreak{}5cc5a237\allowbreak{}a3361f18\allowbreak{}f45fabf1\allowbreak{}4cd23a1c\allowbreak{}679556f6\allowbreak{}5477f9f6\allowbreak{}81915d02; execution/task\_manifest.jsonl SHA‑256 = 5a822e78\allowbreak{}7302a0b3\allowbreak{}bb03a86a\allowbreak{}81b20564\allowbreak{}a44e2636\allowbreak{}731f853f\allowbreak{}9d45ac12\allowbreak{}e4876cc8; analysis/paired\_contingency\_table.csv SHA‑256 = 8fb135ce\allowbreak{}c541cbe0\allowbreak{}b14f1612\allowbreak{}5db33154\allowbreak{}a41460c4\allowbreak{}9958427f\allowbreak{}e8e7f3e1\allowbreak{}f2abfece; analysis/analysis\_log.jsonl SHA‑256 = 261ec80e\allowbreak{}a9343fcc\allowbreak{}bcd0a622\allowbreak{}94f53827\allowbreak{}73770fd5\allowbreak{}9d387dc9\allowbreak{}d749de50\allowbreak{}ca909b3e. The human‑readable verifier codebook (violation‑code \ensuremath{\rightarrow} description) and the full list of artifact paths and hashes are enumerated in the archived execution manifest (execution/execution\_manifest.json); the execution manifest lists the exact verifier codebook artifact path and its SHA‑256 for direct retrieval. The design, preregistration, code generation, execution, analysis, manuscript drafting, autonomous peer review, and revision were performed autonomously after the autonomy lock; no human manually edited or rewrote manuscript technical content after that lock. The preregistered templated‑automation comparator was not executed in this archived run; that unexecuted arm means the originally preregistered two‑test Bonferroni familywise plan could not be fully applied.

\begin{thebibliography}{99}
\bibitem{ref1} doi:10.1145/3718958.3750537
\bibitem{ref2} Sebastian Simon, Alina Mailach, Johannes Dorn, Norbert Siegmund, "A Methodology for Evaluating RAG Systems: A Case Study On Configuration Dependency Validation", 2024, 10.48550/arxiv.2410.08801
\bibitem{ref3} Nguyen Tu, Sukhyun Nam, James Won‐Ki Hong, "Intent‐Based Network Configuration Using Large Language Models", 2024, 10.1002/nem.2313
\bibitem{ref4} Lingzhe Zhang, Tong Jia, Mengxi Jia, Yifan Wu, Aiwei Liu, Yong Yang, Zhonghai Wu, Xuming Hu, Philip S. Yu, Ying Li, "A Survey of AIOps in the Era of Large Language Models", 2025, 10.1145/3746635
\bibitem{ref5} Zhaodong Wang, Samuel Lin, Guanqing Yan, Soudeh Ghorbani, Minlan Yu, Jiawei Zhou, Nathan Hu, Lopa Baruah, Sam Peters, Srikanth Kamath, Jian Yang, Ying Zhang, "Intent-Driven Network Management with Multi-Agent LLMs: The Confucius Framework", 2025, 10.1145/3718958.3750537
\end{thebibliography}

\end{document}
